\documentclass[10pt,a4paper]{amsart}
\usepackage[margin=19mm]{geometry}
\usepackage{amsmath,amssymb,mathtools}
\usepackage[scr=rsfs,cal=euler]{mathalfa}
\usepackage{xfrac}
\usepackage[unicode,hidelinks,bookmarks=false]{hyperref}
\newcommand{\ie}{i.e.\ }
\newcommand{\cf}{cf.\ }
\newcommand{\resp}{respectively\ }
\newcommand{\Subsection}{\subsection}
\providecommand{\nobreakdash}{\nobreak\hbox{-}}
\setlength{\emergencystretch}{2em}
\multlinegap0pt
\usepackage{bm}
\usepackage{bbm}
\usepackage{dsfont}
\usepackage{xspace}
\usepackage{cancel}
\usepackage{mathtools}
\usepackage{cleveref}
\usepackage{enumitem}
\usepackage{amsthm}
\makeatletter
\@ifundefined{theorem}{\newtheorem{theorem}{Theorem}}{}
\@ifundefined{lemma}{\newtheorem{lemma}[theorem]{Lemma}}{}
\@ifundefined{proposition}{\newtheorem{proposition}[theorem]{Proposition}}{}
\renewenvironment{proof}[1]{\par\medskip\noindent\textbf{#1}\enspace\ignorespaces}{\par\medskip}
\makeatother
\newcommand\cA{\mathcal{A}}
\newcommand\cC{\mathcal{C}}
\newcommand\cM{\mathcal{M}}
\newcommand\dR{\mathbb{R}}
\newcommand\dd{\mathrm{d}}
\newcommand{\E}{{\mathbb E}}
\newcommand{\NRM}[1]{{{\left\| #1\right\|}}} % ||1||
\newcommand{\PAR}[1]{{{\left(#1\right)}}} % (1)
\renewcommand{\P}{{\mathbf P}}
\title[Measure estimation on a manifold explored by a diffusion pro]{Measure estimation on a manifold explored by a diffusion process}
\author{Vincent Divol, H\'el\`ene Gu\'erin, Dinh-Toan Nguyen, Viet Chi Tran}
\date{Source: arXiv:2410.11777v2 (CC BY 4.0)}
\begin{document}
\maketitle

\noindent\emph{Source and licence.} Measure estimation on a manifold explored by a diffusion process. Version arXiv:2410.11777v2 (CC BY 4.0), distributed under the Creative Commons Attribution 4.0 International licence, as recorded in the arXiv licence metadata for this exact version and shown on the abstract page https://arxiv.org/abs/2410.11777v2. Full rights notice: licenses/arxiv-2410.11777-CC-BY-4.0.txt.\par\smallskip

\noindent\textit{Editorial scope.} Counted unit: the Lemma whose source label is \texttt{prop:proba-any-distribution} in the article (source0.tex lines 1153--1156, source label \texttt{prop:proba-any-distribution}), delivered together with the article's own complete original proof of it (source0.tex lines 1159--1176, one \texttt{proof} environment of 18 lines). No mathematics is written, repaired or re-derived: statement and proof are the article's own text line for line. Required context retained uncounted: eq:def-K\_h (lines 364--366); lem:properties-Kh (lines 704--712); prop:bernstein (lines 828--831); prop:proba-invariante-unif (lines 878--884). Every definition, display and label of those lines is retained unchanged. Omitted from this package and not redistributed: 1--363, 367--703, 713--827, 832--877, 885--1152, 1157--1158, 1177--2267. Nothing inside a retained excerpt is omitted or altered: every retained body is delivered exactly as the article's lines are written, so no omission receipt of any kind accompanies this package. Citations: the retained excerpts carry no citation command at all, so no bibliography entry of the article is transcribed and no reference is redistributed. Reference adaptation: none was needed and none was made; every reference inside the delivered unit resolves inside it, so no unresolved reference or citation is produced. Printed slips kept literal: no printed word, symbol, hypothesis or number of the counted statement or of its proof was altered. Layout and class: the article's own class is \texttt{article} with packages including geometry, color, array, enumitem, inputenc, amsmath, amssymb, amsthm, booktabs, adjustbox, multirow, microtype, cleveref, babel; the wrapper is the lane's pinned \texttt{amsart} editorial wrapper at 10pt on A4, which restates the theorem-like environments the retained text needs and adds the article's own macro definitions that the retained text needs (\texttt{\textbackslash E}, \texttt{\textbackslash NRM}, \texttt{\textbackslash P}, \texttt{\textbackslash PAR}, \texttt{\textbackslash cA}, \texttt{\textbackslash cC}, \texttt{\textbackslash cM}, \texttt{\textbackslash dR}, \texttt{\textbackslash dd}), with extra wrapper packages (bm, bbm, dsfont, xspace, cancel, mathtools, cleveref, enumitem). The delivered numbering is the unit's own. Assets: no retained span contains a figure environment, a graphics-inclusion command, a PDF-page inclusion, a table or a TikZ picture; no image file is copied into this package, the package declares no asset dependency and no image file is redistributed with it. Licence: version arXiv:2410.11777v2 of this article is distributed under the Creative Commons Attribution 4.0 International licence, as recorded in the arXiv licence metadata for this exact version and shown on the abstract page https://arxiv.org/abs/2410.11777v2. A full rights notice is delivered at licenses/arxiv-2410.11777-CC-BY-4.0.txt. Characters: every retained excerpt is pure ASCII, so the delivered engine, XeLaTeX with its default Latin Modern text font, reports no missing character.
\begin{equation}\label{eq:def-K_h}
K_h(x,y):=\frac{1}{\eta_h(x)}K\PAR{\frac{\NRM{x-y}}{h}},
\end{equation}

\begin{lemma}\label{lem:properties-Kh}
Assume that $K$ is a continuous function on $\dR$ with support in $[0,1]$, such that $\int_{\dR^d}K(\NRM{z})\dd z=1$. Let $h>0$, and consider $K_h$ defined by \eqref{eq:def-K_h} with the renormalizing factor $\eta_h$. Then,
\begin{enumerate}[label=(\roman*)]
\item $h^{-d}\eta_h$ converges to $1$ uniformly on $\cM$;
\item {there exists $h_c>0$, depending on $\cM$ and $K$, such that $\forall h<h_c$, $K_h$ is well defined and} $\forall x\in \cM, \displaystyle{\int_\cM K_h(x,y)\dd y=1}$; 
\item  $\forall h<h_c$, $K_h$ is bounded on $\cM\times \cM$ with  $\NRM{K_h}_\infty\le 2\NRM{K}_\infty h^{-d}$;
\item if furthermore $K$ is Lipschitz continuous, then $\forall h<h_c$, and all $x\in \cM$,  $y\mapsto K_h(x,y)$ is Lipschitz continuous with constant $2\mathrm{Lip}(K)h^{-d-1}$, where $\mathrm{Lip}(K)$ is the Lipschitz constant of $K$.
\end{enumerate}
\end{lemma}

	\begin{proposition}\label{prop:bernstein} 
		Assume that $d\ge 4$. There exist $c_0 $ depending on $\cM$, $K$, $p_{\min}$, $p_{\max}$ and $\kappa_{\min}$, and $h_0$ depending on $\cM$, $K$, $p_{\min}$ and the $\cC^1$-norm of $p$ such that  for any $T>0$ and $y \in \cM$, when $0<h\le h_0$, we have
		$$\P_{\mu}\PAR{ p_{T,h}(y) <p_{\min}/4} \le \exp\PAR{- c_0 T h^{d-2} }.$$
	\end{proposition}

	\begin{proposition}\label{prop:proba-invariante-unif}
 Assume that $d\ge 4$. There exists $h_0$ depending on $\cM$, $K$,  $p_{\min}$ and  the $\cC^1$-norm of $p$ such that  for any $T>0$ and $0<h\le h_0$, we have
\[
 \P_{\mu}\PAR{\inf_{y\in \cM}p_{T,h}(y)<p_{\min}/8 } \le c_1 h^{-d(d+1)} \exp\PAR{- c_0 T h^{d-2} }
\]
for some constant $c_1$ depending only on $p_{\min}$, $K$ and $\cM$, where $c_0$ is the constant of \Cref{prop:bernstein}.
	\end{proposition} 

	\begin{lemma}\label{prop:proba-any-distribution}
 Assume that $d\ge 4$. Let $h_0$, $c_0$ and $c_1$ be  the constants of  \Cref{prop:proba-invariante-unif}. Then, for any $T\ge 2$, if $0<h\le h_0$ and $Th^d>32\|K\|_\infty/p_{\min}$, we have
	\[ \P_{x}\PAR{ E_{T,h}^c }  \le u_\cA c_1h^{-d(d+1)}\exp\PAR{- \frac{c_0}{2}Th^{d-2}}.\]
\end{lemma} 

\begin{proof}{When $T\ge 2$, we have that $1/2\le (T-1)/T\le 1$.
Hence we have the following observation for all $y \in \cM$ and $T\ge 2$:}
\begin{align*}
	p_{T,h}(y)= & \frac{1}{T}\int_0^1 K_h(X_s,y)\dd s +\frac{T-1}{T} \times \frac{1}{T-1}  \int_0^{T-1}  K_h(X_{s+1},y)\dd s \\
	\ge &  - \frac{1}{T} \int_0^1 \| K_h\|_{\infty} \dd s + \frac{{C}}{T-1}  \int_0^{T-1}  K_h(X_{s+1},y)\dd s,
\end{align*}{where $C=1/2$ if $\int_0^{T-1}  K_h(X_{s+1},y)\dd s>0$ or $C=1$ otherwise. }Thus,
	\begin{align*}
		\P_{x}\PAR{ \inf_{ y\in\cM}p_{T,h}(y) <0 } & \le  \P_{x} \PAR{   \inf_{y \in \cM } \PAR{ \frac{{C}}{T-1}  \int_0^{T-1} K_h(X_{s+1},y)\dd s } <\frac{\|K_h\|_{\infty}}{T }}  \\
		 &\le \E_{x} \PAR{\P_{X_1}\PAR{ \inf_{ y\in\cM}p_{(T-1),h}(y) < \frac{2\|K_h\|_{\infty}}{ T }}}\\
&\le u_\cA \P_{\mu}\PAR{ \inf_{ y\in\cM}p_{(T-1),h}(y) < \frac{2\|K_h\|_{\infty}}{  T }},
	\end{align*}
where we used  ultracontractivity at the last step. By Lemma~\ref{lem:properties-Kh}, we have $\|K_h\|_{\infty}\le 2\|K\|_\infty h^{-d}$ for $h$ small enough. Hence, if $Th^d>32\|K\|_\infty/  p_{\min}$, then $2\|K_h\|_{\infty}/T <p_{\min}/8$, and we obtain thanks to \Cref{prop:proba-invariante-unif} that
\begin{align*}
    \P_{x}\PAR{E_{T,h}^c}&\le u_\cA \P_{\mu}\PAR{ \inf_{ y\in\cM}p_{(T-1),h}(y) < p_{\min}/8}\\
&\le u_\cA c_{1}h^{-d(d+1)}\exp(-c_0(T-1)h^{d-2}).
\end{align*}
As $T-1\ge T/2$, for $T\ge 2$, the result holds.
\end{proof}

\end{document}
